% Recovery provenance, 2026-11-20: this is a complete reconstruction from
% the authored proof chunks, the full root audit, and the subsequently
% preserved A.159 proof.  It is a new source identity, not a claim to have
% recovered the lost historical 24481-character TeX byte stream.
% The source-dependent interfaces have been rechecked in the original
% Mnev arXiv:0708.4039v3 during reconstruction.
\Needspace{10\baselineskip}
\subsection{A.147. Parameter regularity and finite PL straightening}

\begin{saresult}{A.147}{Proved under stated hypotheses}{A precise regularization criterion for higher PL comparison}
Fix an o-minimal expansion of the real field. Let \(B\) be a compact
closed PL manifold, put \(N=B\times I\), and let \(k\geq1\).
A class in \(\pi_kP_{\mathrm{def}}(B)\) belongs to the image of
\[
 \pi_kP_{\mathrm{PL}}(B)\longrightarrow\pi_kP_{\mathrm{def}}(B)
 \tag{147.1}
\]
if and only if it has a based definable representative
\(U:N\times P\to N\), \(P=[0,1]^k\), for which a constant \(L<\infty\)
satisfies
\[
 \|U_p(x)-U_q(x)\|_\infty\leq L\|p-q\|_\infty
 \quad(x\in N,\ p,q\in P).
 \tag{147.2}
\]
Here \(U_p\) fixes \(B\times\{0\}\) pointwise, preserves the top end,
and \(U_p=\operatorname{id}\) for \(p\in\partial P\).
Only the forward parameter estimate is assumed.
In particular, real semialgebraic
\(\pi_2P_{\mathrm{PL}}(T^6)\to\pi_2P_{\mathrm{SA}}(T^6)\)
is surjective exactly when each based semialgebraic two-cube is based
semialgebraically homotopic to one satisfying (147.2).
This criterion does not itself assert existence of such representatives.
\end{saresult}

\begin{proof}
All metrics in this proof are maximum norms in fixed finite PL models.
Changes of such models multiply real Lipschitz constants by finite
constants. We give the finite construction, including its relative
boundary conventions; we do not assume a parametric Hauptvermutung.

\textbf{A definable fixed-point fact.}\quad
If \(T:[0,1]^a\to[0,1]^a\) is continuous and definable and has
contraction constant \(c<1\), it has one fixed point. Indeed,
\(\|T(t)-t\|_\infty\) has a minimum \(m\). At a minimizing \(t_0\),
the assumption \(m>0\) would give
\[
 \|T(T(t_0))-T(t_0)\|_\infty\leq cm<m.
\]
Uniqueness follows from the contraction estimate. If \(t,t'\) are
the fixed points for \(T,T'\) with a common contraction bound, then
\[
 \|t-t'\|_\infty\leq
 (1-c)^{-1}\sup_z\|T(z)-T'(z)\|_\infty.
 \tag{147.3}
\]
Consequently a jointly continuous definable family of such maps has
a jointly continuous definable fixed point. Definability follows from
uniqueness; continuity follows from (147.3) and uniform continuity on
the compact variable cube, locally in the other variables. This proof
uses a minimum, not an infinite limiting iteration.

\textbf{The double and the finite Hudson functions.}\quad
For a connected positive-dimensional \(B\), double \(N\) along its two
ends, obtaining the closed PL manifold
\[
 X=N^+\cup N^-\cong B\times S^1.
\]
Write \(A_0,A_1\) for the disjoint seams. Reflect \(U\) to the second
half. The resulting family \(V\) preserves \(N^\pm,A_0,A_1\), fixes
\(A_0\) pointwise, and is the identity on \(\partial P\).
The fixed PL reflection changes (147.2) by only a finite factor.

Choose a finite PL triangulation \(Q\) containing all four marks as
subcomplexes. Choose the preliminary subdivisions needed for the
small-ball separation construction with at most \(k\) support stars;
this is the finite subdivision in Mn\"ev, Proposition~57,
printed pp.~89--90. It depends on \(Q,k\), before a number of time
levels is chosen. Let \(K\) be a sufficiently fine such subdivision,
also refining \(\operatorname{sd}Q\).
Enumerate its vertices \(v_1,\ldots,v_\nu\), and let \(\lambda_i\)
be their continuous PL barycentric functions.
For an integer \(M>0\), use the sequence
\[
 f^{a\nu+i}(y)=\frac aM+
        \frac1M\sum_{j=1}^{i}\lambda_j(y),
 \quad 0\leq a<M,\quad0\leq i\leq\nu .
 \tag{147.4}
\]
The coincident values at successive block ends are identified.
There are \(R=M\nu+1\) values, with
\[
 0=f^0\leq f^1\leq\cdots\leq f^{R-1}=1.
\]
Each consecutive difference is supported in one closed vertex star.
The spatial Lipschitz constants of all \(f^j\) are bounded by
\(M^{-1}\sum_i\operatorname{Lip}(\lambda_i)\).
Choose \(M\) so that this common bound \(c\) satisfies \(cL<1\).
This is legitimate over the real field. The later all-field version
must separately justify a standard bound before choosing this integer.

\textbf{An explicit graph deformation.}\quad
For \(0\leq v\leq1\) and \(0\leq q\leq T_v\), put
\[
 \rho_{y,v}(q)=
   \min_{0\leq j<R}\max\{f^j(y),q-2jv\},
 \qquad T_v=1+2(R-1)v.
 \tag{147.5}
\]
These are finite PL formulas. The values lie in \([0,1]\);
\(\rho_{y,v}(0)=0\), \(\rho_{y,v}(T_v)=1\),
and \(\rho_{y,0}(q)=q\) for \(0\leq q\leq1\).
They have spatial Lipschitz bound \(c\).
For \((x,p,v)\in X\times P\times I\), solve
\[
 t_i=\rho_{V_t(x),v}(T_vp_i),\qquad 1\leq i\leq k,
 \qquad H_{p,v}(x)=V_t(x).
 \tag{147.6}
\]
The right side is a contraction in \(t\) with constant \(cL<1\).
The fixed-point fact supplies unique joint continuous definable
solutions. Moreover
\[
 H_{p,v}^{-1}(y)=
 V^{-1}_{(\rho_{y,v}(T_vp_i))_{i=1}^k}(y).
 \tag{147.7}
\]
For this formula the proposed parameter vector is exactly the vector
required in (147.6); uniqueness verifies both inverse identities.
The original total family inverse is jointly continuous, since its
total map is a continuous bijection of compact spaces. Thus (147.7)
also proves joint continuity of the inverse family.

We have \(H_{p,0}=V_p\). The deformed maps preserve the halves and
seams. On \(A_0\), every \(V_t\) is the identity, so \(H\) is also
the identity. If a coordinate of \(p\) is an endpoint, the same
coordinate of \(t\) is that endpoint. Thus on every parameter boundary
face \(V_t=\operatorname{id}\), and the deformation is based.
No estimate for \(V^{-1}\) in the parameter variable has entered.

\textbf{Finite support labels at the terminal time.}\quad
Set \(D=2R-1\), and at \(v=1\) write \(q_i=Dp_i\).
For each moving band
\(2j\leq q_i\leq2j+1\), \(0\leq j<R-1\), formula (147.5) is
\[
 \rho_{y,1}(q_i)=
 \min\{f^{j+1}(y),\max\{f^j(y),q_i-2j\}\}.
 \tag{147.8}
\]
On the next plateau it equals \(f^{j+1}(y)\).
The last index contributes only the final constant plateau.
Within an interval meeting at most one moving band, the value in
(147.8) is independent of \(q_i\) outside the star supporting
\(f^{j+1}-f^j\).

Cover the finite interval \([0,D]\) by finitely many relatively open
PL intervals, each meeting at most one moving band. Take their
coordinate products and all nonempty intersections. Each is a box.
For such a box \(E_a\), let \(S_a\) be its actual list of moving-band
support stars, one for each active coordinate; repetitions are harmless.
It has at most \(k\) members. Crucially, the label of an intersection
is computed from the bands met by that intersection, rather than
by taking the union of the original box labels. Hence
\[
 E_a\subset E_b\quad\Longrightarrow\quad S_a\subset S_b .
 \tag{147.9}
\]
Formula (147.7) is pointwise independent of \(p\in E_a/D\) outside
the union \(J_a\) of the stars in \(S_a\). These are finite support
data depending only on \(K,M,k\), not on the individual values of \(V\).
They are the open PL parameter topology and monotone support
presheaf used in the finite construction.

Mn\"ev's Proposition~57 replaces configurations of at most \(k\) small
stars by compatible configurations of pairwise disjoint conic PL balls
containing those stars. Its natural transformation is in the covering
order. We retain the entire finite indexing poset, including indices
with repeated labels, and make the inclusions strict by the following
finite version of the conic enlargement used in Proposition~35,
pp.~61--62. Start with maximal indices and enlarge their balls slightly,
keeping the balls in each configuration disjoint and consistent with
the fixed triangulation. Proceed in decreasing rank. At a lower
index each old ball is compactly contained in the already enlarged
higher balls that contain it. A sufficiently small conic enlargement
stays in all of these interiors and retains disjointness within the
lower configuration. The finitely many requirements can be imposed
simultaneously. This gives strict containment for proper indexing
relations. We use this explicit extension of the conic construction;
we do not apply Proposition~35 to a map lacking its strictness
hypothesis.

Write \(\xi_a\) for the resulting configurations. Outside their union,
the inverse family is pointwise constant on the corresponding box.
Continuity gives pointwise constancy on the boundaries as well.
Relative to one fixed parameter in the connected box, the resulting
homeomorphisms fix the complement and the boundaries of the disjoint
balls. They preserve each ball: the possible permutation of finitely
many ball interiors is locally constant in the parameter and starts
at the identity. Thus the inverse image of every ball is constant.
This is the configuration-prismaticity needed below.

\textbf{A split marked coarsening.}\quad
The preceding separation must respect the boundary of \(N\); applying
an unmarked closed-manifold statement and then restricting would not
suffice. In the present application the input coloring is constant:
\(Q(\lambda)=Q\) for every index to which Proposition~36 is applied.
Thus the marked \(Q\) and the requirement \(K\preceq\operatorname{sd}Q\)
were fixed before \(M\) and \(\xi\) were chosen. The rank construction
below produces new subdivision complexes, but it does not require
choosing \(K\) again or making it refine those new complexes.
All input cells already respect \(N^\pm,A_0,A_1\).
Mn\"ev's Proposition~36, pp.~62--63, supplies its conic-star
construction and the incidence property that a consistent ball
centered at \(b\) meets a \(Q(\lambda)\)-cell only when that cell
contains \(b\). A ball meeting a seam must therefore have its center
in that seam and cannot meet both disjoint seams.
Conic-star invariance, together with local flatness of the seam as
a PL subcomplex, identifies the marked pair with a ball and an
equatorial disk. This local-pair conclusion is an inference from
these facts, not an additional literal assertion in Proposition~36.
The two intersections with the halves are its closed half-balls.

In the coarsening construction of Lemma~6, pp.~63--65, replace a
crossing ball's single interior cell by the two half-ball interiors
and the middle-disk interior, retaining their refined boundary cells.
For a ball avoiding the marks use the original single interior cell.
Every new closed cell is a PL ball; its boundary is a union of the
retained cells. The fine complex refines this split complex because
it already respects the marks. This is a split marked coarsening,
not the unchanged inscribed-ball conclusion of Lemma~6.

The decreasing-rank proof is still finite. At a given rank triangulate
each boundary compatibly with its equatorial disk; extend that boundary
triangulation to the complement, refining the original target cells,
the previously constructed higher-rank subdivisions, and the fixed
marks. The interiors of the present balls are left unchanged.
For nested balls, intersection with the same global half or seam
preserves every inclusion, so their split cells have the required
refinement incidence. This is precisely the comparison needed in
the final step of the proof of Lemma~6.

There is also a direct check of prismaticity after splitting. If \(A\)
is a seam or a half, mark preservation gives
\[
 H_p^{-1}(D_0\cap A)=H_p^{-1}(D_0)\cap A
 \tag{147.10}
\]
for every configuration ball \(D_0\). Thus the split interiors have
constant inverse images. Arbitrary chosen boundary subcells have
constant inverse images because the inverse is pointwise constant
on the ball boundary. That last assertion follows from constancy on
the exterior and continuity; the other balls have positive separation.
No assertion about connected components of a split preimage is needed.

\textbf{From the configurations to literal product prisms.}\quad
Apply the finite procedure in the proof of Mn\"ev's Lemma~8,
pp.~67--68, with no auxiliary \(G_i\), with \(Z=\varnothing\), and
with a constant marked initial target coloring.
Here is the assembly, so that its use does not conceal a new
approximation step. Replace the finite open PL parameter topology
by a strict closed PL approximation (Proposition~26). Strengthen it
with the closed-cell topology of a finite triangulation of \(P\).
Proposition~25 realizes the resulting finite order data on an order
complex, simultaneously recording the target coloring and \(\xi\).
On that order complex perform the decreasing-rank split construction
above, obtaining the two natural transformations consisting of a
subdivision and a coarsening. Paste their two mapping-cylinder
colorings as in Step~4 of Lemma~8. Restrict the resulting coloring
to its terminal copy of \(P\).

All choices just listed use only the finite target coloring,
configuration presheaf, and parameter topology. Their validity for
the present definable input follows from the inverse-image
constancies (147.7)--(147.10): the construction never approximates
the input maps or differentiates them. Thus the same proof gives a
finite terminal coloring for this definable homeomorphism family.
If several families have the same data and constancies, the choices
may literally be made once for all of them.

The resulting target prisms are the products \(C\times\sigma\) of
Proposition~4, p.~21. Equation~(11) and Proposition~5, pp.~22--23,
identify their full inverse images with
\[
 H^{-1}(C\times\sigma)=D_C^\sigma\times\sigma,
 \tag{147.11}
\]
where \(D_C^\sigma\) is independent of the parameter on the entire
closed simplex \(\sigma\). These are finite ball complexes.
The spatial cells in the source are at this stage definable balls.
The whole boundary of each product is a union of proper lower
dimensional prisms; on a lower parameter face its spatial
decomposition may be refined. This causes no omission of boundary
data in the induction below. The marks are unions of these prisms.

\textbf{One ambient change makes the spatial cells PL.}\quad
Take the finite family of all source spatial cell closures and all
cells of the fixed target complex \(Q\), including the marks.
A compatible definable triangulation
\(\tau:|K'|\to X\) makes every preimage in this list polyhedral.
Apply marked Hauptvermutung, Shiota, arXiv:1002.1508,
Theorem~2.1, printed p.~5, to the pair of compact polyhedra with
the finite marks \(\tau^{-1}Q_i\) and \(Q_i\).
It gives a PL homeomorphism \(a:|K'|\to X\) carrying those marks
to the corresponding \(Q_i\). Then
\[
 \phi=\tau a^{-1}:X\to X
 \tag{147.12}
\]
preserves every \(Q_i\) setwise, and
\(\phi^{-1}(D)=a\tau^{-1}(D)\) is a subpolyhedron for every source
spatial cell. It is a PL ball: apply the same marked
Hauptvermutung to its definable ball chart and boundary pair.
Thus common conjugation by \(\phi\) makes the source and target
spatial cells PL without losing any fixed mark.

The cell-preserving map \(\phi\) is definably isotopic to the identity,
preserving every \(Q_i\). This follows from the finite cellular
extension construction in the next paragraph, with the parameter
space a point and PL endpoint the identity. It is not an appeal to
an unstated ambient isotopy complement of Hauptvermutung.
Conjugating \(H\) along this isotopy preserves the parameter boundary
identity and the pointwise bottom identity. Call the resulting
terminal family \(W\).

\textbf{Finite endpoint extension by total product coning.}\quad
Construct a PL map \(G\) with the same product-prism correspondence
as \(W\), by increasing total dimension. Prescribe the identity
first on
\[
 E=(A_0\times P)\cup(X\times\partial P).
\]
This is a subcomplex, and identity is the required correspondence
there. For a remaining source prism \(D_0\times\sigma\), all of
\[
 (\partial D_0\times\sigma)\cup(D_0\times\partial\sigma)
       =\partial(D_0\times\sigma)
\]
is already assigned. Make compatible finite boundary subdivisions
so that this prescribed boundary map is PL on the displayed products;
previous maps themselves are unchanged. Choose fixed PL charts of
\(D_0\) and the target spatial ball \(C\) to convex PL balls.
Cone the entire product boundary from centers
\((d_0,b_\sigma)\) and \((c_0,b_\sigma)\), using a common base
barycentre. The extension is a PL homeomorphism. It preserves the
base projection exactly, since corresponding cone segments have
parameter \((1-u)b_\sigma+up\).
This is the full-boundary extension in Mn\"ev's Lemma~3,
proof Steps~2--3, pp.~24--25. Coning only the spatial boundary
with a parameter-dependent formula is not substituted for this
total-product construction.
Finite pasting yields the desired jointly PL \(G\).

\textbf{Relative definable straightening of the remaining cellular map.}\quad
Set \(K=G^{-1}W\). It preserves each source prism and equals the
identity on \(E\). Construct a homotopy \(K\to\operatorname{id}\)
by increasing total prism dimension, relative to \(E\).
For one spatial ball \(C\) and base simplex \(\sigma\), put
\(D_1=\sigma\times I\), where the extra coordinate is homotopy time.
The already assigned face homotopies, initial \(K\), and final identity
give compatible data on
\[
 (\partial C\times D_1)\cup(C\times\partial D_1).
 \tag{147.13}
\]
Use fixed PL ball charts to centered cubes. The prescribed spatial
boundary homeomorphisms extend radially to homeomorphisms \(R_p\)
of the whole spatial cube for every \(p\in D_1\).
For \(p\in\partial D_1\), divide the prescribed whole-ball map by
\(R_p\), obtaining a homeomorphism \(Q_p\) fixing the spatial boundary.
Identify the parameter ball \(D_1\) with a centered cube. Write
\(p=ru\), \(r=\|p\|_\infty\), \(u\in\partial D_1\), when \(r>0\).
The formula
\[
 A_r(Q_u)(x)=
 \begin{cases}
   rQ_u(x/r),&\|x\|_\infty\leq r,\quad r>0,\\
   x,&\|x\|_\infty\geq r,
 \end{cases}
 \qquad A_0=\operatorname{id}
 \tag{147.14}
\]
is a homeomorphism fixing the spatial boundary; the two formulas
agree where they overlap. Its displacement is at most \(2r\),
so it is jointly continuous at the parameter center even though
\(u\) is undefined there. The family \(R_pA_r(Q_u)\), with value
\(R_0\) at the center, extends all of (147.13). It is definable,
and all inverse families are continuous by the same formulas
or the compact total-map argument.
Finite pasting proves the relative homotopy.

This extension argument also establishes the earlier isotopy of
\(\phi\) to the identity, by applying it successively to the cells
of \(Q\). No spatial or parameter smoothness is used.
Multiplying the constructed homotopy by \(G\) joins \(W\) to \(G\)
while preserving the marks and \(E\). Concatenate the graph
deformation and the conjugation homotopy, and restrict from \(X\)
to \(N^+\). This proves the sufficient direction of (147.1).

\textbf{Converse and degenerate cases.}\quad
A jointly PL map on a compact real polyhedron is globally Lipschitz
in fixed PL coordinates, hence satisfies (147.2). A class in the
comparison image therefore has such a representative.
For disconnected \(B\), basedness and connectedness of \(P\) force
each component to remain fixed; apply the construction on finitely
many positive-dimensional components and paste.
For a point component, the concordance is an increasing homeomorphism
of the interval fixing its ends. The convex homotopy
\((1-t)U_p(x)+tx\) is again such a homeomorphism, is based, and ends
at the PL identity. Empty bases have only their empty maps.
These observations complete all cases and the stated named
\(T^6,\pi_2\) criterion.
\end{proof}

\textbf{Source scope.}\quad
The source is Mn\"ev, \emph{Combinatorial fiber bundles and fragmentation
of a fiberwise PL-homeomorphism}, arXiv:0708.4039v3:
Proposition~4, equation~(11), Proposition~5 and Lemma~3,
pp.~21--25; Propositions~35--36 and Lemma~6, pp.~61--65;
Lemma~8, pp.~67--68; Proposition~57, pp.~89--90.
The original statements concern PL input. The proof above supplies
the definable graph deformation, the marked split variant, and the
finite relative cellular straightening required to use the finite
combinatorics here. Marked Hauptvermutung is used only after a
finite spatial-cell system has been obtained; it is not asserted
to preserve an arbitrary parameter projection.
